%% toc: P1 7 7.1,7.2,7.3
\section{Proofs for ambiguity, restricted quantification, uniqueness, and scope in translation}
\label{sec:annex-p01-w02-p01-022-translation}

Notation and definitions are those of
\cref{sec:p01-w02-p01-022-translation}. For a structure \(\mathcal M\), an
assignment \(s\), a formula \(\theta\), a variable \(x\) and
\(a\in\dom\mathcal M\), the phrase ``\(\theta\) holds at \(a\)'' abbreviates
\(\mathcal M,s[x\mapsto a]\entails\theta\) when \(\mathcal M\), \(s\) and
\(x\) are fixed by the context. Satisfaction clauses are those of
\cref{sec:fol-semantics-02}; \(\theta\leftrightarrow\theta'\) abbreviates
\((\theta\to\theta')\land(\theta'\to\theta)\), and repeated \(\land\) and
\(\lor\) associate to the right
\cite[ch.~``Propositional Syntax and Semantics'', definition ``Defined
connectives'' and the association convention preceding
it]{FrozenFormalLogicFormalizing}. For a sentence \(\sigma\),
\(\mathcal M,s\entails\sigma\) iff \(\mathcal M\entails\sigma\), by the
definition of truth of a sentence in \cref{sec:fol-semantics-02}.

\begin{annexitem}{annex:p01-w02-p01-022-readings}{thm:p01-w02-p01-022-readings}{Validity under alternative readings}
\begin{proof}
\emph{Clause 1, premises.} Assume \(\varphi_1,\varphi_2\) are
\(\Pi\)-equivalent and \(\Gamma\cup\Pi\cup\{\varphi_1\}\entails\psi\). Let
\(\mathcal M,s\) satisfy every member of \(\Gamma\cup\Pi\cup\{\varphi_2\}\).
Every member of \(\Pi\) is a sentence, so \(\mathcal M\) is a model of
\(\Pi\); by \(\Pi\)-equivalence, \(\mathcal M,s\entails\varphi_1\). Hence
\(\mathcal M,s\) satisfies \(\Gamma\cup\Pi\cup\{\varphi_1\}\), and so
\(\mathcal M,s\entails\psi\). This proves
\(\Gamma\cup\Pi\cup\{\varphi_2\}\entails\psi\). Exchanging the indices gives
the converse.

\emph{Clause 1, conclusions.} Assume \(\psi_1,\psi_2\) are
\(\Pi\)-equivalent and \(\Gamma\cup\Pi\entails\psi_1\). If \(\mathcal M,s\)
satisfies \(\Gamma\cup\Pi\), then \(\mathcal M,s\entails\psi_1\), and
\(\mathcal M\) is a model of \(\Pi\), so \(\mathcal M,s\entails\psi_2\).
Exchanging the indices gives the converse.

\emph{An auxiliary fact.} For \(k\geq1\),
\(\mathcal M,s\entails\varphi_1\lor\cdots\lor\varphi_k\) iff
\(\mathcal M,s\entails\varphi_i\) for some \(i\leq k\). For \(k=1\) the
iterated disjunction is \(\varphi_1\). For \(k>1\) it is
\(\varphi_1\lor(\varphi_2\lor\cdots\lor\varphi_k)\); by the disjunction
clause it is satisfied iff \(\varphi_1\) is or
\(\varphi_2\lor\cdots\lor\varphi_k\) is, and by the induction hypothesis on
\(k\) the latter holds iff some \(\varphi_i\) with \(2\leq i\leq k\) is
satisfied. In the same way, using the conjunction clause,
\(\mathcal M,s\entails\psi_1\land\cdots\land\psi_m\) iff
\(\mathcal M,s\entails\psi_j\) for every \(j\leq m\).

\emph{Clause 2.} Assume \(\Gamma\cup\Pi\cup\{\varphi_i\}\entails\psi\) for
every \(i\). Let \(\mathcal M,s\) satisfy
\(\Gamma\cup\Pi\cup\{\varphi_1\lor\cdots\lor\varphi_k\}\). By the auxiliary
fact some \(\varphi_i\) is satisfied, so \(\mathcal M,s\) satisfies
\(\Gamma\cup\Pi\cup\{\varphi_i\}\) and hence \(\psi\). Conversely, assume the
disjunctive consequence and fix \(i\). If \(\mathcal M,s\) satisfies
\(\Gamma\cup\Pi\cup\{\varphi_i\}\), it satisfies the disjunction by the
auxiliary fact, hence \(\psi\).

\emph{Clause 3.} Assume \(\Gamma\cup\Pi\entails\psi_j\) for every \(j\). If
\(\mathcal M,s\) satisfies \(\Gamma\cup\Pi\), it satisfies every \(\psi_j\),
hence the conjunction by the auxiliary fact. Conversely, if the conjunction
is a consequence and \(\mathcal M,s\) satisfies \(\Gamma\cup\Pi\), then
\(\mathcal M,s\) satisfies the conjunction and hence each \(\psi_j\).
\end{proof}
\end{annexitem}

\begin{annexitem}{annex:p01-w02-p01-022-negation-scope}{thm:p01-w02-p01-022-negation-scope}{Negation scope over a restricted universal}
\begin{proof}
Both formulas are sentences, so assignments play no role.

\emph{Clause 1.} Let \(\mathcal M_1\) have domain \(\{0,1\}\),
\(S^{\mathcal M_1}=\{0,1\}\) and \(P^{\mathcal M_1}=\{0\}\). At \(1\),
\(S(x)\) holds and \(P(x)\) fails, so \(S(x)\to P(x)\) fails at \(1\),
\(\forall x\,(S(x)\to P(x))\) is false and \(\rho_2\) is true. At \(0\),
\(S(x)\) holds and \(\neg P(x)\) fails, so \(\rho_1\) is false. Hence
\(\rho_2\not\entails\rho_1\). Let \(\mathcal M_2\) have domain \(\{0\}\) and
\(S^{\mathcal M_2}=P^{\mathcal M_2}=\varnothing\). At \(0\) the antecedent
\(S(x)\) fails, so both \(S(x)\to\neg P(x)\) and \(S(x)\to P(x)\) hold at
\(0\). Thus \(\rho_1\) is true and \(\forall x\,(S(x)\to P(x))\) is true, so
\(\rho_2\) is false. Hence \(\rho_1\not\entails\rho_2\).

\emph{Clause 2.} Let \(\mathcal M\) satisfy \(\rho_1\) and
\(\exists x\,S(x)\). Choose \(a\in S^{\mathcal M}\). By \(\rho_1\),
\(S(x)\to\neg P(x)\) holds at \(a\), and \(S(x)\) holds at \(a\), so
\(a\notin P^{\mathcal M}\). Then \(S(x)\to P(x)\) fails at \(a\), so
\(\forall x\,(S(x)\to P(x))\) is false in \(\mathcal M\) and \(\rho_2\) is
true.
\end{proof}
\end{annexitem}

\begin{annexitem}{annex:p01-w02-p01-022-indefinite}{thm:p01-w02-p01-022-indefinite}{Indefinites in conditional antecedents}
\begin{proof}
Fix \(\mathcal M\) and \(s\). For every \(b\in\dom\mathcal M\), the
assignments \(s[y\mapsto b]\) and \(s\) agree on every variable other than
\(y\), and \(y\notin\FV(\chi)\); by formula coincidence
(\cref{lem:fol-coincidence}),
\(\mathcal M,s[y\mapsto b]\entails\chi\) iff \(\mathcal M,s\entails\chi\).
By the universal and conditional clauses,
\(\mathcal M,s\entails\forall y\,(\varphi\to\chi)\) iff for every \(b\):
\(\mathcal M,s[y\mapsto b]\not\entails\varphi\) or
\(\mathcal M,s[y\mapsto b]\entails\chi\); by the coincidence just noted,
iff for every \(b\): \(\mathcal M,s[y\mapsto b]\not\entails\varphi\) or
\(\mathcal M,s\entails\chi\).

If \(\mathcal M,s\entails\chi\), this condition holds, and
\(\exists y\,\varphi\to\chi\) holds by the conditional clause. If
\(\mathcal M,s\not\entails\chi\), the condition says that no \(b\)
satisfies \(\varphi\) at \(s[y\mapsto b]\), that is,
\(\mathcal M,s\not\entails\exists y\,\varphi\), which by the conditional
clause is exactly what \(\exists y\,\varphi\to\chi\) requires when its
consequent fails. In both cases the two formulas have the same truth value
at \(\mathcal M,s\).
\end{proof}
\end{annexitem}

\begin{annexitem}{annex:p01-w02-p01-022-relativization}{thm:p01-w02-p01-022-relativization}{Semantics of restricted quantifiers}
\begin{proof}
Write ``at \(a\)'' for satisfaction at \(\mathcal M,s[x\mapsto a]\); so
\(a\in A\) iff \(\alpha\) holds at \(a\).

\emph{Clause 1.} By the universal and conditional clauses,
\(\mathcal M,s\entails\forall x\,(\alpha\to\varphi)\) iff for every
\(a\in\dom\mathcal M\), \(\alpha\) fails at \(a\) or \(\varphi\) holds at
\(a\). For \(a\notin A\) the first alternative holds. For \(a\in A\) it
fails, so the condition at \(a\) is that \(\varphi\) holds at \(a\). Hence
the formula is satisfied iff \(\varphi\) holds at every \(a\in A\).

\emph{Clause 2.} By the existential and conjunction clauses,
\(\mathcal M,s\entails\exists x\,(\alpha\land\varphi)\) iff some \(a\) has
\(\alpha\) and \(\varphi\) at \(a\), that is, iff some \(a\in A\) has
\(\varphi\) at \(a\).

\emph{Clause 3.} By the negation clause and clause 1,
\(\mathcal M,s\entails\neg\forall x\,(\alpha\to\varphi)\) iff not every
\(a\in A\) has \(\varphi\) at \(a\), iff some \(a\in A\) has \(\neg\varphi\)
at \(a\), iff (clause 2 applied to \(\neg\varphi\))
\(\mathcal M,s\entails\exists x\,(\alpha\land\neg\varphi)\). Likewise
\(\mathcal M,s\entails\neg\exists x\,(\alpha\land\varphi)\) iff no
\(a\in A\) has \(\varphi\) at \(a\) (clause 2), iff every \(a\in A\) has
\(\neg\varphi\) at \(a\), iff (clause 1 applied to \(\neg\varphi\))
\(\mathcal M,s\entails\forall x\,(\alpha\to\neg\varphi)\). The set \(A\) is
the same on both sides because \(\alpha\), \(x\), \(\mathcal M\) and \(s\)
are. Since \(\mathcal M\) and \(s\) were arbitrary, the formulas are
logically equivalent.

\emph{Clause 4.} If \(A=\varnothing\), the condition of clause 1 holds
vacuously and the condition of clause 2 fails.
\end{proof}
\end{annexitem}

\begin{annexitem}{annex:p01-w02-p01-022-mismatch}{thm:p01-w02-p01-022-mismatch}{The conditional and conjunction mismatches}
\begin{proof}
Fix \(\mathcal M,s\) and write ``at \(a\)'' as before.

\emph{Clause 1.} \(\mathcal M,s\entails\forall x\,(\alpha\land\varphi)\) iff
\(\alpha\) and \(\varphi\) hold at every \(a\). The right-hand side is
satisfied iff \(\alpha\) holds at every \(a\) and, for every \(a\), if
\(\alpha\) holds at \(a\) then \(\varphi\) does. If the left side holds, the
first conjunct of the right side holds, and the second holds because
\(\varphi\) holds at every \(a\). If the right side holds, then for every
\(a\), \(\alpha\) holds at \(a\) by the first conjunct and then \(\varphi\)
holds at \(a\) by the second.

\emph{Clause 2.} Suppose \(\mathcal M,s\entails\exists x\,(\alpha\to
\varphi)\), with \(\alpha\to\varphi\) holding at \(a\). If \(\alpha\) fails
at \(a\), then \(\neg\alpha\) holds at \(a\) and \(\exists x\,\neg\alpha\) is
satisfied. Otherwise \(\alpha\) holds at \(a\), so \(\varphi\) holds at
\(a\), and \(\exists x\,(\alpha\land\varphi)\) is satisfied. Conversely, if
\(\neg\alpha\) holds at some \(a\), then \(\alpha\to\varphi\) holds at \(a\);
and if \(\alpha\land\varphi\) holds at some \(a\), then \(\varphi\) holds at
\(a\) and so does \(\alpha\to\varphi\). Either way
\(\exists x\,(\alpha\to\varphi)\) is satisfied.

\emph{Clauses 3 and 4.} The two entailments follow from clauses 1 and 2 with
\(\alpha=A(x)\) and \(\varphi=Q(x)\): the universal conjunction implies its
second conjunct, and the restricted existential implies the disjunction of
which it is a disjunct. For the converses, let \(\mathcal M\) have domain
\(\{0,1\}\), \(A^{\mathcal M}=\{0\}\) and \(Q^{\mathcal M}=\{0\}\). Then
\(A(x)\to Q(x)\) holds at \(0\) (both parts hold) and at \(1\) (the antecedent
fails), so \(\forall x\,(A(x)\to Q(x))\) is true; but \(A(x)\) fails at
\(1\), so \(\forall x\,(A(x)\land Q(x))\) is false. Let \(\mathcal M'\) have
domain \(\{0\}\) and \(A^{\mathcal M'}=Q^{\mathcal M'}=\varnothing\). Then
\(A(x)\to Q(x)\) holds at \(0\), so \(\exists x\,(A(x)\to Q(x))\) is true,
while \(A(x)\land Q(x)\) holds nowhere, so \(\exists x\,(A(x)\land Q(x))\) is
false.

\emph{Clauses 5 and 6.} In \(\mathcal M'\), \(\forall x\,(A(x)\to Q(x))\) is true
by clause 4 of \cref{thm:p01-w02-p01-022-relativization} and
\(\exists x\,(A(x)\land Q(x))\) is false, which proves the non-entailment.
For the entailment, let \(\mathcal M\) satisfy
\(\forall x\,(A(x)\to Q(x))\) and \(\exists x\,A(x)\), and choose
\(a\in A^{\mathcal M}\). The universal premise gives \(A(x)\to Q(x)\) at
\(a\), so \(a\in Q^{\mathcal M}\), and \(A(x)\land Q(x)\) holds at \(a\).
\end{proof}
\end{annexitem}

\begin{annexitem}{annex:p01-w02-p01-022-uniqueness}{thm:p01-w02-p01-022-uniqueness}{Counting characterization of uniqueness}
\begin{proof}
\emph{Step 0: evaluation at a pair.} Let \(a,b\in\dom\mathcal M\) and put
\(t=s[x\mapsto a][y\mapsto b]\). Since \(x\ne y\), \(t(x)=a\), \(t(y)=b\) and
\(t(z)=s(z)\) for every other variable \(z\).
\begin{enumerate}
  \item[(i)] \(\mathcal M,t\entails\varphi\) iff \(a\in S\). Indeed \(t\)
        and \(s[x\mapsto a]\) agree on every variable except possibly \(y\),
        and \(y\notin\FV(\varphi)\) because \(y\) does not occur in
        \(\varphi\); formula coincidence (\cref{lem:fol-coincidence}) gives
        \(\mathcal M,t\entails\varphi\) iff
        \(\mathcal M,s[x\mapsto a]\entails\varphi\), that is, iff \(a\in S\).
  \item[(ii)] \(\mathcal M,t\entails\varphi[y/x]\) iff \(b\in S\). By the
        substitution lemma (\cref{thm:fol-substitution}) with the term \(y\),
        \(\mathcal M,t\entails\varphi[y/x]\) iff
        \(\mathcal M,t[x\mapsto t(y)]\entails\varphi\), and \(t(y)=b\). The
        assignment \(t[x\mapsto b]\) sends \(x\) to \(b\), \(y\) to \(b\) and
        every other variable \(z\) to \(s(z)\); it agrees with
        \(s[x\mapsto b]\) on every variable except possibly \(y\). By formula
        coincidence as in (i), \(\mathcal M,t[x\mapsto b]\entails\varphi\)
        iff \(\mathcal M,s[x\mapsto b]\entails\varphi\), that is, iff
        \(b\in S\).
  \item[(iii)] By the equality clause, \(\mathcal M,t\entails x=y\) iff
        \(a=b\), and \(\mathcal M,t\entails y=x\) iff \(b=a\).
\end{enumerate}

\emph{Clause 1.} By the two universal clauses and the conditional and
conjunction clauses, \(\mathcal M,s\entails\exists^{\leq1}x\,\varphi\) iff
for all \(a,b\in\dom\mathcal M\), if \(\mathcal M,t\entails\varphi\) and
\(\mathcal M,t\entails\varphi[y/x]\) then \(\mathcal M,t\entails x=y\),
where \(t=s[x\mapsto a][y\mapsto b]\). By Step 0 this says: for all \(a,b\),
if \(a\in S\) and \(b\in S\) then \(a=b\). That is exactly the statement that
\(S\) has at most one element.

\emph{Clause 2.} By the existential and conjunction clauses,
\(\mathcal M,s\entails\exists!x\,\varphi\) iff there is \(a\) with
\(\mathcal M,s[x\mapsto a]\entails\varphi\) and
\(\mathcal M,s[x\mapsto a]\entails\forall y\,(\varphi[y/x]\to y=x)\). The
first conjunct says \(a\in S\). By the universal and conditional clauses and
Step 0 (ii) and (iii), the second says: for every \(b\), if \(b\in S\) then
\(b=a\), that is, \(S\subseteq\{a\}\). Hence the formula is satisfied iff
\(S=\{a\}\) for some \(a\), that is, iff \(S\) has exactly one element.

\emph{Clause 3.} By the existential clause,
\(\mathcal M,s\entails\exists x\,\varphi\) iff \(S\ne\varnothing\). With
clause 1, \(\exists x\,\varphi\land\exists^{\leq1}x\,\varphi\) is satisfied
iff \(S\) is nonempty and has at most one element, iff \(S\) has exactly
one element, iff (clause 2) \(\exists!x\,\varphi\) is satisfied. For the
second formula, \(\mathcal M,s\entails\exists x\,\forall y\,
(\varphi[y/x]\leftrightarrow y=x)\) iff there is \(a\) such that for every
\(b\), \(\mathcal M,s[x\mapsto a][y\mapsto b]\) satisfies both
\(\varphi[y/x]\to y=x\) and \(y=x\to\varphi[y/x]\). By Step 0 this says:
for every \(b\), \(b\in S\) iff \(b=a\); that is, \(S=\{a\}\). By clause 2
this is the condition for \(\exists!x\,\varphi\). As \(\mathcal M\) and \(s\)
were arbitrary, the three formulas are logically equivalent.

\emph{Independence of the comparison variable.} The conditions in clauses 1
and 2 mention only \(S=\varphi^{\mathcal M,s}_x\), which is defined without
reference to \(y\). Two comparison variables therefore give formulas with
the same truth value at every \(\mathcal M,s\), that is, logically
equivalent formulas.
\end{proof}
\end{annexitem}

\begin{annexitem}{annex:p01-w02-p01-022-freshness}{thm:p01-w02-p01-022-freshness}{The comparison variable must be fresh}
\begin{proof}
The substitution \((x=y)[y/x]\) acts on the atomic formula by replacing the
term \(x\) by \(y\) and leaving the term \(y\) unchanged, so it yields
\(y=y\), and \(\chi\) is as displayed.

\emph{The solution set.} Since \(x\ne y\), \(s[x\mapsto a](y)=s(y)\), so by
the equality clause \(\mathcal M,s[x\mapsto a]\entails x=y\) iff
\(a=s(y)\). Hence \(\varphi^{\mathcal M,s}_x=\{s(y)\}\).

\emph{The formula \(\chi\).} Suppose \(\mathcal M,s\entails\chi\). Then there
is \(a\) with \(\mathcal M,s[x\mapsto a]\entails\forall y\,(y=y\to y=x)\),
so for every \(b\in\dom\mathcal M\), the assignment
\(s[x\mapsto a][y\mapsto b]\) satisfies \(y=y\to y=x\). It satisfies
\(y=y\), since \(b=b\); hence it satisfies \(y=x\), so \(b=a\). Thus every
element of \(\dom\mathcal M\) equals \(a\), contradicting the assumption that
\(\mathcal M\) has two distinct elements. Therefore
\(\mathcal M,s\not\entails\chi\).
\end{proof}
\end{annexitem}

\begin{annexitem}{annex:p01-w02-p01-022-exactly-one}{thm:p01-w02-p01-022-exactly-one}{Exactly one P is Q, and the P is Q}
\begin{proof}
\emph{\(E_1\) and the abbreviation.} The variable \(y\) does not occur in
\(P(x)\land Q(x)\), and \((P(x)\land Q(x))[y/x]\) is \(P(y)\land Q(y)\). So
\(\exists!x\,(P(x)\land Q(x))\) is
\(\exists x\,((P(x)\land Q(x))\land\forall y\,((P(y)\land Q(y))\to y=x))\),
which differs from \(E_1\) only in the association of the conjunctions. By
the conjunction clause, both associations are satisfied at an assignment
iff all three conjuncts are, so the two formulas are logically equivalent.
The solution set of \(P(x)\land Q(x)\) in \(x\) is
\(P^{\mathcal M}\cap Q^{\mathcal M}\); by
\cref{thm:p01-w02-p01-022-uniqueness}, clause 2, \(\mathcal M\entails E_1\)
iff \(P^{\mathcal M}\cap Q^{\mathcal M}\) has exactly one element.

\emph{\(E_2\entails E_1\).} Let \(\mathcal M\entails E_2\) with witness
\(a\). Then \(a\in P^{\mathcal M}\) and \(a\in Q^{\mathcal M}\), and for
every \(b\), \(s[x\mapsto a][y\mapsto b]\) satisfies \(P(y)\to y=x\), so
\(b\in P^{\mathcal M}\) implies \(b=a\). Thus \(P^{\mathcal M}=\{a\}\) and
\(P^{\mathcal M}\cap Q^{\mathcal M}=\{a\}\), so \(\mathcal M\entails E_1\).

\emph{\(E_1\not\entails E_2\).} Let \(\mathcal M\) have domain \(\{0,1\}\),
\(P^{\mathcal M}=\{0,1\}\) and \(Q^{\mathcal M}=\{0\}\). Then
\(P^{\mathcal M}\cap Q^{\mathcal M}=\{0\}\), so \(\mathcal M\entails E_1\).
A witness \(a\) for \(E_2\) would need \(b\in P^{\mathcal M}\Rightarrow b=a\)
for both \(b=0\) and \(b=1\), which is impossible; so
\(\mathcal M\not\entails E_2\).

\emph{\(\exists x\,(P(x)\land Q(x))\not\entails E_1\).} Let \(\mathcal M\)
have domain \(\{0,1\}\) and \(P^{\mathcal M}=Q^{\mathcal M}=\{0,1\}\). Then
\(P(x)\land Q(x)\) holds at \(0\), so the existential sentence is true, while
\(P^{\mathcal M}\cap Q^{\mathcal M}\) has two elements, so \(E_1\) is false.
\end{proof}
\end{annexitem}

\begin{annexitem}{annex:p01-w02-p01-022-description-scope}{thm:p01-w02-p01-022-description-scope}{Truth conditions and scope of eliminated descriptions}
\begin{proof}
\emph{Clause 1.} By the existential and conjunction clauses, the formula
\(\mathsf{the}_x(\varphi,\chi)\) is satisfied at \(\mathcal M,s\) iff there
is \(a\) such
that, at \(s[x\mapsto a]\), the formulas \(\varphi\),
\(\forall y\,(\varphi[y/x]\to y=x)\) and \(\chi\) are satisfied. As in the
proof of clause 2 of \cref{thm:p01-w02-p01-022-uniqueness} (which uses only
that \(y\ne x\) does not occur in \(\varphi\)), the first two hold iff
\(a\in S\) and \(S\subseteq\{a\}\), that is, iff \(S=\{a\}\). The third is
\(\mathcal M,s[x\mapsto a]\entails\chi\). The variable \(y\) may occur in
\(\chi\); this does not matter, because \(\chi\) is evaluated at
\(s[x\mapsto a]\), outside the scope of \(\forall y\).

\emph{Clause 2.} By clause 1, satisfaction of
\(\mathsf{the}_x(\varphi,\chi)\) requires \(S=\{a\}\) for some \(a\), which
fails when \(S\) is empty or has two or more elements, whatever \(\chi\) is.

\emph{Clause 3.} Let \(\mathcal M,s\entails\mathsf{the}_x(\varphi,\neg\chi)\).
By clause 1, \(S=\{a\}\) and \(\mathcal M,s[x\mapsto a]\not\entails\chi\).
If \(\mathcal M,s\entails\mathsf{the}_x(\varphi,\chi)\), clause 1 would give
\(S=\{a'\}\) with \(\mathcal M,s[x\mapsto a']\entails\chi\); then
\(a'=a\), a contradiction. Hence
\(\mathcal M,s\entails\neg\,\mathsf{the}_x(\varphi,\chi)\).

\emph{Clause 4.} Let \(\mathcal M,s\) satisfy \(\exists!x\,\varphi\) and
\(\neg\,\mathsf{the}_x(\varphi,\chi)\). By clause 2 of
\cref{thm:p01-w02-p01-022-uniqueness}, \(S=\{a\}\) for some \(a\). If
\(\mathcal M,s[x\mapsto a]\entails\chi\), clause 1 would give
\(\mathcal M,s\entails\mathsf{the}_x(\varphi,\chi)\), which is excluded. So
\(\mathcal M,s[x\mapsto a]\entails\neg\chi\), and clause 1 gives
\(\mathcal M,s\entails\mathsf{the}_x(\varphi,\neg\chi)\).

\emph{Clause 5.} If \(S=\varnothing\), clause 2 applied to \(\chi\) and to
\(\neg\chi\) shows that \(\mathsf{the}_x(\varphi,\chi)\) and
\(\mathsf{the}_x(\varphi,\neg\chi)\) are both false at \(\mathcal M,s\); so
the negation of the first is true and the second is false.
\end{proof}
\end{annexitem}

\begin{annexitem}{annex:p01-w02-p01-022-worked}{thm:p01-w02-p01-022-worked}{Verdicts for the worked example}
\begin{proof}
All three formulas are sentences. Write \(\varphi\) for
\(D(y)\land B(y,a)\); the description variable is \(y\) and the comparison
variable \(z\) does not occur in \(\varphi\). For a structure \(\mathcal M\)
and any assignment \(s\), the solution set of \(\varphi\) in \(y\) is
\(S=\{c\in D^{\mathcal M}:(c,a^{\mathcal M})\in B^{\mathcal M}\}\).

\emph{Validity.} Let \(\mathcal M\) satisfy \(\mathrm{P}_1\) and
\(\mathrm{P}_2\), and fix an assignment \(s\). By clause 1 of
\cref{thm:p01-w02-p01-022-description-scope} (with \(y\) as description
variable and \(z\) as comparison variable), \(S=\{c\}\) for some \(c\) and
\(\mathcal M,s[y\mapsto c]\entails\exists x\,(R(x)\land H(x,y))\). Choose
\(r\) with \(r\in R^{\mathcal M}\) and \((r,c)\in H^{\mathcal M}\). Since
\(c\in S\), \(c\in D^{\mathcal M}\). By \(\mathrm{P}_1\), the assignment
\(s[x\mapsto r][y\mapsto c]\) satisfies
\((R(x)\land D(y)\land H(x,y))\to G(x,y)\); its antecedent holds, so
\((r,c)\in G^{\mathcal M}\). Hence \(s[y\mapsto c][x\mapsto r]\) satisfies
\(R(x)\land G(x,y)\), and \(\mathcal M,s[y\mapsto c]\entails
\exists x\,(R(x)\land G(x,y))\). Since \(S=\{c\}\), clause 1 of the same
theorem gives \(\mathcal M,s\entails\mathrm{C}\).

\emph{Invalidity with \(\mathrm{P}_2'\).} Let \(\mathcal M\) have domain
\(\{r,p_1,p_2,n\}\) of four distinct elements, with \(a^{\mathcal M}=n\),
\(R^{\mathcal M}=\{r\}\), \(D^{\mathcal M}=\{p_1,p_2\}\),
\(B^{\mathcal M}=\{(p_1,n),(p_2,n)\}\) and
\(H^{\mathcal M}=G^{\mathcal M}=\{(r,p_1),(r,p_2)\}\). The pairs \((u,v)\)
with \(u\in R^{\mathcal M}\), \(v\in D^{\mathcal M}\) and
\((u,v)\in H^{\mathcal M}\) are \((r,p_1)\) and \((r,p_2)\), both in
\(G^{\mathcal M}\); at every other pair the antecedent of \(\mathrm{P}_1\)
fails. So \(\mathcal M\entails\mathrm{P}_1\). The element \(p_1\) is in
\(D^{\mathcal M}\), \((p_1,n)\in B^{\mathcal M}\), \(r\in R^{\mathcal M}\)
and \((r,p_1)\in H^{\mathcal M}\), so \(\mathcal M\entails\mathrm{P}_2'\).
But \(S=\{p_1,p_2\}\) has two elements, so clause 2 of
\cref{thm:p01-w02-p01-022-description-scope} gives
\(\mathcal M\not\entails\mathrm{C}\).
\end{proof}
\end{annexitem}
